\documentclass{amsart}
\usepackage{amssymb,mathtools}
\usepackage[hidelinks]{hyperref}
\hypersetup{pdftitle={The irrationality exponent of small-height binary lacunary series is 2},pdfauthor={Ruiyi Zhang}}

\newtheorem{theorem}{Theorem}
\newtheorem{lemma}{Lemma}[section]
\numberwithin{equation}{section}
\newcommand{\Q}{\mathbb Q}
\newcommand{\Z}{\mathbb Z}
\newcommand{\R}{\mathbb R}
\newcommand{\sgn}{\operatorname{sgn}}

\begin{document}

\title[Small-height binary lacunary series]{The irrationality exponent of small-height binary lacunary series is 2}
\author{Ruiyi Zhang}
\address{Independent Researcher, Columbia, Maryland, USA}
\email{ryzfsu@gmail.com}
\subjclass[2020]{Primary 11J82; Secondary 11J70, 11B83}
\keywords{irrationality exponent, lacunary series, rational height, Hankel determinant, Pad\'e approximation}

\begin{abstract}
Let $N\ge2$ be an integer and let $\lambda_j$ be nonzero rational
numbers with $\log H(\lambda_j)=o(2^j)$, where $H(u/v)=\max(|u|,v)$
for $u/v$ in lowest terms with $v>0$. We prove that $\sum_{j\ge0}\lambda_jN^{-2^j}$ is irrational and
has irrationality exponent $2$. The proof uses Pad\'e approximants to
the tails $\sum_{r\ge0}\lambda_{k+r}z^{2^r}$. The two Hankel families
they need satisfy parity recursions, so they are nonzero for all nonzero
weights.
After the finite prefix of the series is cleared, the approximants of
order $n$ give integer linear forms with height
$N^{(n+1/2)2^k}$ and error $N^{-(n+1/2)2^k}$, up to factors
$e^{o(2^k)}$. Only finitely many orders are used for each target
exponent. The height condition cannot be weakened to
$\log H(\lambda_j)\le C2^j$ for any constant $C>0$: for every $C$ some
such weights give exponent greater than $2$, and for larger $C$
positive integer weights can give a rational number or a Liouville
number.
\end{abstract}

\maketitle

\section{Introduction}

For an irrational real number $\xi$, its irrationality exponent is
\[
\mu(\xi)=\sup\left\{\nu\in\R:
\left|\xi-\frac pq\right|<q^{-\nu}
\text{ for infinitely many }(p,q)\in\Z\times\Z_{>0}\right\}.
\]
The pigeonhole principle gives $\mu(\xi)\ge2$ (see the end of Section~5).
Write $L=\log N$.
For a nonzero rational number in lowest terms, put
\begin{equation}\label{eq:rationalheight}
H(u/v)=\max(|u|,v),\qquad u\in\Z\setminus\{0\},\quad
v\in\Z_{>0},\quad \gcd(u,v)=1.
\end{equation}
We show that equality holds for binary lacunary series with rational
weights of sufficiently small height.

\begin{theorem}\label{thm:main}
Let $N\in\Z$, $N\ge2$, and let $(\lambda_j)_{j\ge0}$ be a sequence
in $\Q\setminus\{0\}$ such that $\log H(\lambda_j)=o(2^j)$.
Then the absolutely convergent series
\begin{equation}\label{eq:x}
x=\sum_{j=0}^{\infty}\lambda_jN^{-2^j}
 =\frac{\lambda_0}{N}+\frac{\lambda_1}{N^2}
  +\frac{\lambda_2}{N^4}+\frac{\lambda_3}{N^8}+\cdots
\end{equation}
is irrational and satisfies $\mu(x)=2$.
\end{theorem}

Geometric weights $\lambda_j=\lambda^j$, with $\lambda\in\Q\setminus\{0\}$,
satisfy the hypothesis because $\log H(\lambda^j)=O(j)$.
Every sequence drawn in arbitrary order from a finite subset of
$\Q\setminus\{0\}$ also satisfies the hypothesis, since its heights are bounded.

Duverney's criterion \cite[Corollary~3.4]{Du01} already gives
irrationality under this height condition. For $w_j=N^{2^j}$, its
rationality exception would require
\[
\lambda_j\lambda_{j+2}=w_j\lambda_{j+1}^2
\]
eventually. Here
$\log|\lambda_j\lambda_{j+2}/\lambda_{j+1}^2|=o(2^j)$, whereas
$\log w_j=2^jL$, so the exception cannot occur.
Duverney \cite[Theorem~3.2]{Du01} also studies finite exponent bounds.
For related work, see his later paper \cite{Du24}.

Duverney, Kurosawa and Shiokawa \cite[Theorem~1]{DKS} consider signed
series $\sum\epsilon_j/x_j$, where $x_j>1$ are rational and
$x_{j+1}\ge x_j^2$ eventually. Their denominator condition is
\begin{equation}\label{eq:DKS}
\delta_1=\operatorname{den}(x_1),\qquad
\delta_{j+1}=\delta_j^2\operatorname{den}(x_{j+1}/x_j^2),\qquad
\log\delta_{j+1}=o(\log x_j),
\end{equation}
where $\operatorname{den}$ is the reduced positive denominator.
They conclude
$\mu(\sum\epsilon_j/x_j)=\limsup_j\log x_{j+1}/\log x_j$.
For $x_j=N^{2^j}/(t|\lambda_j|)$ with a fixed rational $t>0$,
the ratio is $t|\lambda_j|^2/|\lambda_{j+1}|$; the accumulated
condition in \eqref{eq:DKS} can fail even for a finite set of weights.

Amou and Bugeaud \cite[Lemma~3]{AB} give an earlier criterion for
$\sum B_j/A_j$, with strictly increasing positive integers $A_j$
and nonzero integers $B_j$. Their hypotheses include ratios
$S_j=(A_{j+1}/|B_{j+1}|)/(A_j/|B_j|)^2\ge2$ eventually and
$\prod_{r=0}^j d_r^{2^{j-r}}=A_j^{o(1)}$, where
$d_j=\operatorname{den}(S_j)$ and $d_0=|B_1|$; small heights alone
do not imply this cumulative arithmetic condition.

For unit weights, Shallit \cite[Theorems~8--9]{Sh} proved bounded
partial quotients when $N\ge3$. For $N=2$, he states the corresponding
result for $2\sum_{j\ge0}2^{-2^j}$
\cite[Theorem~11 and following example]{Sh}.
He also states a recursion for weights $(-1)^j$ when $N\ge3$
\cite[Theorem~12]{Sh}.

For nonzero integral geometric weights, Bugeaud, Han, Wen and Yao
\cite[Theorem~2.1]{BHWY} give exponent $2$:
$g(z)=\sum_{r\ge0}\lambda^rz^{2^r-1}$ has integer coefficients,
satisfies $g(z)=1+\lambda zg(z^2)$, and has nonzero Hankel
determinants of every order by Lemma~\ref{lem:hankel}.
The multiplier $\lambda z$ is nonzero on the orbit of $1/N$.
The same theorem covers periodic rational weights. If the period is
$p\ge1$ and $D_{\mathrm{per}}\in\Z_{>0}$ clears the denominators, then
$f(z)=D_{\mathrm{per}}\sum_{j\ge0}\lambda_jz^{2^j-1}\in\Z[[z]]$
converges in the unit disk and satisfies
$f(z)=P(z)+z^{2^p-1}f(z^{2^p})$ with $P\in\Z[z]$.
Lemma~\ref{lem:hankel} gives every Hankel order, and $z^{2^p-1}$
does not vanish on the orbit of $1/N$. Since
$x=f(1/N)/(D_{\mathrm{per}}N)$, \cite[Theorem~2.1]{BHWY} gives
$\mu(x)=2$.
Su \cite[Theorem~7.3]{Su} gives a formal continued fraction for
$\sum_{j\ge0}c_jz^{2^j-1}$ with $c_0=1$ and nonzero $c_j$ in any
field.

Rajchert \cite{Ra} and Su \cite{Su} study related questions on irrationality exponents and continued fractions of Mahler-type series. None of our results depends on these works; apart from standard facts, every step is proved here.

For geometric weights, Mahler's method also gives transcendence
\cite[Corollary~2]{Ni90}. The series
$F(z)=\sum_{r\ge0}\lambda^rz^{2^r}$ has radius $1$ and Hadamard
gaps, so it is transcendental over $\mathbb C(z)$.
It satisfies $F(z)=z+\lambda F(z^2)$; the matrix of the system for
$(1,F(z))^{\mathsf T}$ and its inverse have polynomial entries.
The orbit of $1/N$ therefore contains no singular point.

Our proof uses moving tails and their Hankel recursions. After clearing
the finite prefix, adjacent Pad\'e approximants give integer pairs
with nonzero determinant. For each target exponent the finite set of
orders is fixed first, so all coefficient costs are $o(2^k)$ at tail
shift $k$. Section~6 shows that the height scale is sharp.

Throughout the proof write
\begin{equation}\label{eq:weights}
\begin{gathered}
\lambda_j=\frac{u_j}{v_j},\qquad
u_j\in\Z\setminus\{0\},\quad v_j\in\Z_{>0},\quad
\gcd(u_j,v_j)=1,\\
h_j=\log H(\lambda_j).
\end{gathered}
\end{equation}
In particular,
\begin{equation}\label{eq:weightbounds}
e^{-h_j}\le|\lambda_j|\le e^{h_j}.
\end{equation}
The upper bound and $h_j=o(2^j)$ give absolute convergence of
\eqref{eq:x}. The moving tails are
\begin{equation}\label{eq:F}
F_k(z)=\sum_{r\ge0}\lambda_{k+r}z^{2^r},\qquad
F_k(z)=\lambda_kz+F_{k+1}(z^2).
\end{equation}
For each fixed $k$, the series converges absolutely for $|z|<1$.

\section{Hankel determinants}

Put
\begin{equation}\label{eq:a}
g_k(z)=\frac{F_k(z)}z=\sum_{i\ge0}a_i^{(k)}z^i.
\end{equation}
By \eqref{eq:F},
\begin{equation}\label{eq:coeff}
a_0^{(k)}=\lambda_k,\qquad a_{2i}^{(k)}=0\quad(i\ge1),\qquad
a_{2i+1}^{(k)}=a_i^{(k+1)}\quad(i\ge0).
\end{equation}
Define
\begin{equation}\label{eq:HK}
\begin{aligned}
H_n^{(k)}&=\det(a_{i+j}^{(k)})_{0\le i,j<n},& H_0^{(k)}&=1,\\
K_n^{(k)}&=\det(a_{i+j+1}^{(k)})_{0\le i,j<n},& K_0^{(k)}&=1.
\end{aligned}
\end{equation}

\begin{lemma}\label{lem:hankel}
For all shifts $k\ge0$, the determinants satisfy
\begin{align}
H_{2m}^{(k)}&=(-1)^m(H_m^{(k+1)})^2 &&(m\ge1),\label{eq:Heven}\\
H_{2m+1}^{(k)}&=(-1)^m\lambda_k(K_m^{(k+1)})^2
&&(m\ge0),\label{eq:Hodd}\\
K_n^{(k)}&=H_{\lceil n/2\rceil}^{(k+1)}K_{\lfloor n/2\rfloor}^{(k+1)}
&&(n\ge1).\label{eq:Krec}
\end{align}
In particular, $H_n^{(k)}K_n^{(k)}\ne0$ for every $n,k\ge0$.
\end{lemma}

\begin{proof}
Ordering rows and columns by parity makes the two permutation signs cancel.
For $H_{2m}^{(k)}$, the matrix becomes
\[
\begin{pmatrix}\lambda_k E_{00}&A\\ A^{\mathsf T}&0\end{pmatrix},
\qquad A=(a_{r+s}^{(k+1)})_{0\le r,s<m},
\]
where $E_{00}$ has its sole nonzero entry $1$ in the upper left corner.
Exchanging the two blocks of $m$ columns contributes
$(-1)^{m^2}=(-1)^m$ and leaves diagonal blocks $A$ and
$A^{\mathsf T}$. This proves \eqref{eq:Heven}.

For $H_{2m+1}^{(k)}$, separate index $0$, the remaining even indices,
and the odd indices. The case $m=0$ gives $H_1^{(k)}=\lambda_k$.
For $m\ge1$, put
\[
\boldsymbol a=(a_0^{(k+1)},\ldots,a_{m-1}^{(k+1)})^{\mathsf T},\qquad
G_m=(a_{r+s+1}^{(k+1)})_{0\le r,s<m}.
\]
The matrix is
\begin{equation}\label{eq:oddblock}
\begin{pmatrix}
\lambda_k&0&\boldsymbol a^{\mathsf T}\\
0&0&G_m\\
\boldsymbol a&G_m^{\mathsf T}&0
\end{pmatrix}.
\end{equation}
Since $\lambda_k\ne0$, pivoting on the first entry gives the Schur
complement
\[
\begin{pmatrix}0&G_m\\G_m^{\mathsf T}&-\lambda_k^{-1}
\boldsymbol a\boldsymbol a^{\mathsf T}\end{pmatrix}.
\]
For square blocks of size $m$, column exchange gives
\begin{equation}\label{eq:blockidentity}
\det\begin{pmatrix}0&B\\C&X\end{pmatrix}
=(-1)^m\det B\det C.
\end{equation}
The determinant of \eqref{eq:oddblock} is
$(-1)^m\lambda_k(\det G_m)^2$, proving \eqref{eq:Hodd}.

For $K_n^{(k)}$, the mixed parity blocks vanish by \eqref{eq:coeff}.
The diagonal blocks are
\[
(a_{r+s}^{(k+1)})_{0\le r,s<\lceil n/2\rceil},\qquad
(a_{r+s+1}^{(k+1)})_{0\le r,s<\lfloor n/2\rfloor}.
\]
This proves \eqref{eq:Krec}. Strong induction on $n$, simultaneously
for all $k$, starts from
\[
H_0^{(k)}=K_0^{(k)}=1,\qquad H_1^{(k)}=\lambda_k,\qquad
K_1^{(k)}=\lambda_{k+1}.
\]
Every remaining factor has smaller size, so all determinants are nonzero.
\end{proof}

Only offsets $0$ and $1$ are needed; other Hankel determinants may
vanish, since $a_2^{(k)}=0$. The weights may have either sign.
For geometric weights, $g_k=\lambda^kg_0$; writing $H_n=H_n^{(0)}$
and $K_n=K_n^{(0)}$, the recursions specialize to
\[
H_{2m}=(-1)^m\lambda^{2m}H_m^2,\qquad
H_{2m+1}=(-1)^m\lambda^{2m}K_m^2,\qquad
K_n=\lambda^nH_{\lceil n/2\rceil}K_{\lfloor n/2\rfloor}.
\]

\section{Pad\'e approximants to the moving tails}

\begin{lemma}\label{lem:pade}
For each $n\ge1$ and $k\ge0$, there exist a positive integer $d_{n,k}$
and polynomials $U_{n,k},V_{n,k}\in\Z[z]$ with
\begin{equation}\label{eq:degrees}
\deg U_{n,k},\deg V_{n,k}\le n,\qquad V_{n,k}(0)=d_{n,k}>0,
\end{equation}
such that
\begin{equation}\label{eq:remainder}
E_{n,k}(z):=V_{n,k}(z)F_k(z)-U_{n,k}(z)
=c_{n,k}z^{2n+1}+O(z^{2n+2}),\qquad
c_{n,k}\in\Z\setminus\{0\}.
\end{equation}
They can be chosen so that, for each fixed finite set of orders,
all their coefficients, $d_{n,k}$ and $|c_{n,k}|$ are at most
$\exp(o(2^k))$, uniformly over that set as $k\to\infty$.
They also satisfy the exact identity
\begin{equation}\label{eq:adjacent}
U_{n+1,k}V_{n,k}-U_{n,k}V_{n+1,k}
 =\Omega_{n,k}z^{2n+1},\qquad
\Omega_{n,k}=d_{n+1,k}c_{n,k}\in\Z\setminus\{0\}.
\end{equation}
\end{lemma}

\begin{proof}
Fix $n,k$. By Lemma~\ref{lem:hankel}, the system
\begin{equation}\label{eq:system}
\sum_{i=0}^{n-1}a_{i+j}^{(k)}b_i=-a_{n+j}^{(k)}
\qquad(0\le j<n)
\end{equation}
has a unique rational solution. For a power series $W$, the notation
$[W]_{<n}$ retains only powers below $n$. Put
\begin{gather}
p_{n,k}(T)=T^n+\sum_{i=0}^{n-1}b_iT^i,\qquad
Q_{n,k}(z)=z^np_{n,k}(1/z),\label{eq:pQ}\\
P_{n,k}(z)=z[Q_{n,k}(z)g_k(z)]_{<n}.\label{eq:Ptrunc}
\end{gather}
Both polynomials have degree at most $n$, and $Q_{n,k}(0)=1$.
Here $[z^\ell]W$ denotes the coefficient of $z^\ell$ in $W$.

Define the linear functional $\mathcal L_k(T^j)=a_j^{(k)}$.
For $n\le\ell<2n$, \eqref{eq:system} gives
\[
[z^\ell](Q_{n,k}g_k)=\mathcal L_k(T^{\ell-n}p_{n,k})=0.
\]
At $\ell=2n$, the bordered determinant formula gives
\begin{equation}\label{eq:schur}
\begin{split}
\mathcal L_k(T^np_{n,k})
 &=a_{2n}^{(k)}-\boldsymbol h^{\mathsf T}
       (a_{i+j}^{(k)})_{i,j<n}^{-1}\boldsymbol h
 =\frac{H_{n+1}^{(k)}}{H_n^{(k)}},\\
\boldsymbol h&=(a_n^{(k)},\ldots,a_{2n-1}^{(k)})^{\mathsf T}.
\end{split}
\end{equation}
Hence
\begin{equation}\label{eq:normalizederror}
Q_{n,k}F_k-P_{n,k}
 =\frac{H_{n+1}^{(k)}}{H_n^{(k)}}z^{2n+1}+O(z^{2n+2}).
\end{equation}

We choose the integerizer explicitly, because its size enters the
height. Set
\begin{equation}\label{eq:localwindow}
\ell_n=\lfloor\log_2(2n+1)\rfloor,\qquad
\mathcal D_{n,k}=\prod_{r=0}^{\ell_n}v_{k+r},\qquad
\Sigma_{n,k}=\sum_{r=0}^{\ell_n}h_{k+r}.
\end{equation}
Every $\mathcal D_{n,k}a_i^{(k)}$ with $0\le i\le2n$ is an
integer. Let
\begin{equation}\label{eq:integerizer}
\mathfrak D_{n,k}
 =\det(\mathcal D_{n,k}a_{i+j}^{(k)})_{0\le i,j<n},\qquad
 d_{n,k}=\mathcal D_{n,k}|\mathfrak D_{n,k}|.
\end{equation}
The determinant $\mathfrak D_{n,k}$ is a nonzero integer.
Cramer's rule in \eqref{eq:system} shows that every coefficient of
$d_{n,k}Q_{n,k}$ is an integer multiple of $\mathcal D_{n,k}$.
Consequently
\begin{equation}\label{eq:integerpolys}
U_{n,k}=d_{n,k}P_{n,k},\qquad
V_{n,k}=d_{n,k}Q_{n,k}
\end{equation}
are integer polynomials. The coefficient of degree $2n+1$ in
$V_{n,k}F_k-U_{n,k}$ is also an integer: all coefficients of $F_k$
needed at that degree have denominator dividing $\mathcal D_{n,k}$.
By \eqref{eq:normalizederror}, it is
\begin{equation}\label{eq:c}
c_{n,k}=d_{n,k}\frac{H_{n+1}^{(k)}}{H_n^{(k)}}\ne0.
\end{equation}
In particular, $|c_{n,k}|\ge1$.

Next we bound the sizes in Lemma~\ref{lem:pade}.
Write $\Sigma=\Sigma_{n,k}$. Then
\[
\mathcal D_{n,k}\le e^{\Sigma},\qquad
|\mathcal D_{n,k}a_i^{(k)}|\le e^{2\Sigma}
\quad(0\le i\le2n).
\]
Expansion of the determinant, and of the integer Cramer numerators,
gives
\begin{equation}\label{eq:coefficientcost}
\begin{aligned}
1\le|\mathfrak D_{n,k}|&\le n!e^{2n\Sigma},\\
1\le d_{n,k}&\le n!e^{(2n+1)\Sigma},\\
\max\left\{\max_i|[z^i]U_{n,k}|,
\max_i|[z^i]V_{n,k}|,|c_{n,k}|\right\}
 &\le(n+1)!e^{(2n+2)\Sigma}.
\end{aligned}
\end{equation}
For the last line, a coefficient of $U_{n,k}$ or the leading error
is a convolution of at most $n+1$ terms; a coefficient of $F_k$
contributes at most $e^{\Sigma}$. The division by
$\mathfrak D_{n,k}$ cancels when the Cramer solution is multiplied by
$d_{n,k}$.
Since $n$ is fixed and $h_j=o(2^j)$, $\Sigma_{n,k}=o(2^k)$,
uniformly over finitely many orders. The polynomials, their integerizer
and $c_{n,k}$ depend only on the weights in \eqref{eq:localwindow};
the remainder $E_{n,k}$ depends on the whole tail.

Finally, the polynomial on the left of \eqref{eq:adjacent} has degree
at most $2n+1$. Substituting $U_{j,k}=V_{j,k}F_k-E_{j,k}$ shows that
it has no terms below degree $2n+1$, and that the coefficient at that
degree is $d_{n+1,k}c_{n,k}\ne0$. Thus it is exactly the monomial
in \eqref{eq:adjacent}.
\end{proof}

We evaluate at a point that moves with $k$.
Put
\begin{equation}\label{eq:movingpoint}
\Delta=2^k,\qquad y=N^{-\Delta},\qquad
\varepsilon_k=\sup_{j\ge k}\frac{h_j}{2^j}.
\end{equation}
The hypothesis gives $\varepsilon_k\to0$. Moreover,
\begin{equation}\label{eq:prefixheight}
\sum_{j<k}h_j=o(2^k).
\end{equation}
Given $\rho>0$, choose $J$ such that $h_j\le\rho2^j$ for
$j\ge J$. The fixed sum before $J$, divided by $2^k$, tends to zero,
and the remaining sum divided by $2^k$ is at most $\rho$.

\begin{lemma}\label{lem:local}
For any fixed finite set of orders, uniformly over that set as
$k\to\infty$,
\begin{align}
Q_{n,k}(y)&=1+O\bigl(\exp(-L\Delta+o(\Delta))\bigr),\label{eq:Vbounds}\\
E_{n,k}(y)&=c_{n,k}y^{2n+1}
 \left(1+O\bigl(\exp(-L\Delta+o(\Delta))\bigr)\right).
 \label{eq:Ebounds}
\end{align}
In particular, $V_{n,k}(y)>0$ and $E_{n,k}(y)\ne0$ for all
sufficiently large $k$, and
\begin{equation}\label{eq:errorlog}
\log|E_{n,k}(y)|=-(2n+1)L\Delta+o(\Delta).
\end{equation}
\end{lemma}

\begin{proof}
The integer Cramer numerators used above have absolute value at most
$n!\exp(2n\Sigma_{n,k})$, and their denominator is a nonzero
integer. Thus the nonconstant coefficients of $Q_{n,k}$ have size
$\exp(o(\Delta))$. Since $y=\exp(-L\Delta)$, their finite sum
proves \eqref{eq:Vbounds}. Also $V_{n,k}=d_{n,k}Q_{n,k}$ with
$d_{n,k}>0$, proving eventual positivity.

Fix $n$, and split $F_k$ after the terms with $r\le\ell_n$.
The first omitted exponent is
\[
t_n=2^{\ell_n+1}>2n+1.
\]
For every $r\ge0$, the definition of $\varepsilon_k$ gives
\[
|\lambda_{k+r}|y^{2^r}
 \le\exp\bigl(-(L-\varepsilon_k)\Delta2^r\bigr).
\]
For all sufficiently large $k$, $L-\varepsilon_k\ge L/2$ and
\begin{align}
\sum_{r\ge\ell_n+1}|\lambda_{k+r}|y^{2^r}
&\le\sum_{s\ge0}\exp\bigl(-(L-\varepsilon_k)\Delta t_n2^s\bigr)
 \notag\\
&\le2\exp\bigl(-(L-\varepsilon_k)\Delta t_n\bigr)
 =\exp(-Lt_n\Delta+o(\Delta)).\label{eq:infinitebound}
\end{align}
The second inequality uses $2^s\ge s+1$ and, for large $k$, the
geometric ratio $\exp(-(L-\varepsilon_k)\Delta t_n)\le1/2$.

Multiplying the finite part by $V_{n,k}$ and subtracting
$U_{n,k}$ gives a polynomial whose leading term is
$c_{n,k}z^{2n+1}$. All its subsequent terms have degree at least
$2n+2$ and coefficients of size $\exp(o(\Delta))$ by
\eqref{eq:coefficientcost}. At $y$, their sum is therefore at most
$\exp(-(2n+2)L\Delta+o(\Delta))$.
The omitted part, after multiplication by $V_{n,k}(y)$, is bounded
by \eqref{eq:infinitebound} with an additional factor
$\exp(o(\Delta))$. Since $t_n\ge2n+2$ and $|c_{n,k}|\ge1$,
both parts, relative to $c_{n,k}y^{2n+1}$, have absolute value at
most $\exp(-L\Delta+o(\Delta))$.
This proves \eqref{eq:Ebounds}, and the factor in parentheses is
eventually nonzero. Finally, $0\le\log|c_{n,k}|=o(\Delta)$ proves
\eqref{eq:errorlog}; there are finitely many orders, so all bounds
are uniform.
\end{proof}

\section{Moving-tail iteration and integerization}

Fix $k\ge1$, so that $N^{\Delta/2}$ below is an integer, and retain
$\Delta=2^k$, $y=N^{-\Delta}$.
Repeatedly splitting \eqref{eq:F} gives the exact identity
\begin{equation}\label{eq:iteration}
x=S_k+F_k(y),\qquad
S_k=\sum_{j=0}^{k-1}\frac{u_j}{v_jN^{2^j}}.
\end{equation}
Define
\begin{equation}\label{eq:prefixintegerizer}
C_k=\operatorname{lcm}(v_0,\ldots,v_{k-1}),\qquad
T_k=C_kN^{\Delta/2},\qquad J_k=T_kS_k.
\end{equation}
All three are integers: for $j<k$, $v_j$ divides $C_k$ and
$2^j\le\Delta/2$. Equations \eqref{eq:prefixheight} and
\eqref{eq:prefixintegerizer} give
\begin{equation}\label{eq:prefixcost}
0\le\log C_k\le\sum_{j<k}h_j=o(\Delta),\qquad
\log T_k=\tfrac12 L\Delta+o(\Delta).
\end{equation}

Put
\begin{align}
\widetilde U_{n,k}&=N^{n\Delta}U_{n,k}(y),&
\widetilde V_{n,k}&=N^{n\Delta}V_{n,k}(y),\label{eq:AB}\\
B_{n,k}&=T_k\widetilde V_{n,k},&
A_{n,k}&=J_k\widetilde V_{n,k}+T_k\widetilde U_{n,k}.\notag
\end{align}
The degree bounds give integer evaluations and hence integer pairs
$A_{n,k},B_{n,k}$.
For all sufficiently large $k$ in a fixed finite order set,
$B_{n,k}>0$ and
\begin{align}
\frac{A_{n,k}}{B_{n,k}}
 &=S_k+\frac{U_{n,k}(y)}{V_{n,k}(y)},\label{eq:quotient}\\
B_{n,k}x-A_{n,k}
 &=T_kN^{n\Delta}E_{n,k}(y)\ne0.\label{eq:integererror}
\end{align}

Put
\begin{equation}\label{eq:scale}
\tau_{n,k}=(n+\tfrac12)L\Delta.
\end{equation}
Write the relative error in \eqref{eq:Ebounds} as $r_{n,k}=o(1)$.
Then
\begin{align}
B_{n,k}
 &=C_kd_{n,k}e^{\tau_{n,k}}Q_{n,k}(y),\label{eq:exactheight}\\
B_{n,k}x-A_{n,k}
 &=C_kc_{n,k}e^{-\tau_{n,k}}(1+r_{n,k}).\label{eq:exacterror}
\end{align}
Every additional logarithm in these two expressions is
$o(\Delta)$, uniformly over a fixed finite set of orders. Hence
\begin{align}
\log B_{n,k}&=\tau_{n,k}+o(\Delta),\label{eq:height}\\
\log|B_{n,k}x-A_{n,k}|&=-\tau_{n,k}+o(\Delta),\label{eq:error}\\
\log\left|x-\frac{A_{n,k}}{B_{n,k}}\right|
 &=-2\tau_{n,k}+o(\Delta).\label{eq:rationalerror}
\end{align}
The individual factors have logarithmic size
\[
\begin{array}{c|c}
\text{quantity}&\text{logarithmic size at fixed }n\\ \hline
T_k&\frac12L\Delta+o(\Delta)\\
N^{n\Delta}&nL\Delta\\
d_{n,k},\ |c_{n,k}|&o(\Delta)\\
V_{n,k}(y)&o(\Delta)\\
|E_{n,k}(y)|&-(2n+1)L\Delta+o(\Delta).
\end{array}
\]
The prefix costs $\tfrac12L\Delta$, and the factor $z$ in
$F_k=zg_k$ raises the error order to $2n+1$; the two halves balance,
since $2(n+\tfrac12)=2n+1$. Also
$|A_{n,k}|\le|x|B_{n,k}+|B_{n,k}x-A_{n,k}|$, so $A_{n,k}$
obeys the same height bound.

The adjacent identity \eqref{eq:adjacent} gives
\begin{equation}\label{eq:integerdet}
A_{n+1,k}B_{n,k}-A_{n,k}B_{n+1,k}
=C_k^2N^{\Delta}\Omega_{n,k}\ne0.
\end{equation}
The two terms containing $J_k$ cancel, leaving
\begin{align*}
&T_k^2N^{(2n+1)\Delta}
 \bigl(U_{n+1,k}(y)V_{n,k}(y)-U_{n,k}(y)V_{n+1,k}(y)\bigr)\\
&\qquad=C_k^2N^{(2n+2)\Delta}\Omega_{n,k}y^{2n+1}
 =C_k^2N^{\Delta}\Omega_{n,k}.
\end{align*}
Hence, for every $(p,q)\in\Z\times\Z_{>0}$, the two integers
\begin{equation}\label{eq:sourceintegers}
B_{n,k}p-qA_{n,k},\qquad B_{n+1,k}p-qA_{n+1,k}
\end{equation}
cannot both vanish.

\section{Proof of Theorem~\ref{thm:main}}

First fix $n=1$ and let $k\to\infty$.
By \eqref{eq:integererror} and \eqref{eq:error}, the forms
$B_{1,k}x-A_{1,k}$ are eventually nonzero and tend to zero.
If $x=a/b$ were rational, with $a\in\Z$ and $b\in\Z_{>0}$, then
\[
b(B_{1,k}x-A_{1,k})=B_{1,k}a-bA_{1,k}
\]
would be a nonzero integer tending to zero, a contradiction.
This proves irrationality.

For the upper bound we use a fixed finite set of orders.
Fix an integer $M\ge1$ and use primary orders $M\le n\le2M$,
each with companion $n+1$. The whole order set is therefore
\begin{equation}\label{eq:bank}
\mathcal B=\{M,M+1,\ldots,2M+1\}.
\end{equation}
Put
\begin{equation}\label{eq:alpha}
\alpha_M=1+\frac1{M+1/2}.
\end{equation}
Order the primary scales $\tau_{n,k}$ by increasing $k$, and then
increasing $n$. This is their numerical order, because the largest
scale in one block is $(2M+\tfrac12)L2^k$, and the first in the
next is $(2M+1)L2^k$. Within a block, successive ratios are at most
\[
\frac{M+3/2}{M+1/2}=\alpha_M.
\]
At a block transition the ratio is
\[
\frac{2(M+1/2)}{2M+1/2}<\alpha_M.
\]
The companion scale also satisfies
$\tau_{n+1,k}\le\alpha_M\tau_{n,k}$.

Fix $0<\eta<1$. Since the order set $\mathcal B$ is finite, \eqref{eq:height} and
\eqref{eq:error} hold beyond one common threshold. For each
sufficiently large primary scale $\tau=\tau_{n,k}$, its two rows
$i=n,n+1$ satisfy
\begin{equation}\label{eq:bankbounds}
0<B_{i,k}\le\exp((\alpha_M+\eta)\tau),\qquad
|B_{i,k}x-A_{i,k}|\le\exp(-(1-\eta)\tau).
\end{equation}
The companion has a larger scale, which only improves its error.
Since $\tau$ is a fixed multiple of $\Delta$ for each order in the
order set $\mathcal B$, the $o(\Delta)$ terms are absorbed once $M$ and $\eta$ are fixed.

Let $q$ be any sufficiently large positive integer. Choose the first
primary scale for which
\begin{equation}\label{eq:selection}
\tau\ge\frac{\log(4q)}{1-\eta}.
\end{equation}
For sufficiently large $q$ it has a predecessor beyond the common
threshold. The scale-ratio bound therefore gives
\begin{equation}\label{eq:taurange}
\frac{\log(4q)}{1-\eta}\le\tau
 \le\frac{\alpha_M\log(4q)}{1-\eta}.
\end{equation}
Set
\begin{equation}\label{eq:beta}
\beta=\frac{\alpha_M(\alpha_M+\eta)}{1-\eta}.
\end{equation}
For both selected rows, \eqref{eq:bankbounds}--\eqref{eq:taurange}
imply
\begin{equation}\label{eq:qbounds}
B_{i,k}\le(4q)^\beta,\qquad
q|B_{i,k}x-A_{i,k}|\le\frac14.
\end{equation}

Suppose $p\in\Z$ and
\begin{equation}\label{eq:toogood}
\left|x-\frac pq\right|
 <4^{-(\beta+1)}q^{-(1+\beta)}.
\end{equation}
For $i=n,n+1$, the integer $I_{i,k}=B_{i,k}p-qA_{i,k}$ satisfies
\begin{align}
|I_{i,k}|
&\le B_{i,k}q\left|\frac pq-x\right|
      +q|B_{i,k}x-A_{i,k}|\notag\\
&<(4q)^\beta q\,4^{-(\beta+1)}q^{-(1+\beta)}+\frac14
 =\frac12.\label{eq:contradiction}
\end{align}
Both integers would vanish, contrary to \eqref{eq:integerdet}.
Hence the reverse of \eqref{eq:toogood} holds for every integer $p$
and every sufficiently large $q$, whether or not $p/q$ is reduced or
equals a selected approximant.

Let $\nu>2$. Since $\alpha_M\to1$, choose
$M$ such that $1+\alpha_M^2<\nu$, and then choose $\eta>0$ small
enough that $1+\beta<\nu$. With these choices fixed, the proved
lower bound
\[
\left|x-\frac pq\right|\ge4^{-(\beta+1)}q^{-(1+\beta)}
\]
is larger than $q^{-\nu}$ for all sufficiently large $q$.
Only finitely many pairs with the remaining denominators can satisfy
$|x-p/q|<q^{-\nu}$. Thus $\mu(x)\le\nu$, and letting $\nu$
decrease to $2$ gives $\mu(x)\le2$.

The choices have the order
\begin{equation}\label{eq:quantifiers}
\bigl((\lambda_j)_{j\ge0},N\bigr)
\ \longrightarrow\ \nu>2
\ \longrightarrow\ M,\eta,\mathcal B
\ \longrightarrow\ k_0,q_0
\ \longrightarrow\ (p,q)
\ \longrightarrow\ k,n.
\end{equation}
The thresholds $k_0,q_0$ depend on the sequence and its little-$o$
modulus.

For the reverse bound, fix an integer $R\ge1$ and partition $[0,1)$
into $R$ intervals of length $1/R$.
Two fractional parts among $\{jx\}$, $0\le j\le R$, lie in one interval.
Their difference gives integers $p,q$ with
\[
1\le q\le R,\qquad |qx-p|<1/R,\qquad
\left|x-\frac pq\right|<\frac1{qR}\le\frac1{q^2}.
\]
Since $x$ is irrational, for each fixed $D$ the finitely many
distances $\|qx\|$, $1\le q\le D$, have a positive minimum $m_D>0$,
where $\|t\|$ denotes distance from $t$ to the nearest integer.
For $R>1/m_D$, the inequality $|qx-p|<1/R$ forces $q>D$.
These denominators are unbounded as $R\to\infty$, and
$\mu(x)\ge2$.\qed

\section{Sharpness of the height scale}

The little-$o$ condition cannot be replaced by a bound
$\log H(\lambda_j)\le C_*2^j$ with any constant $C_*>0$. Below,
small coefficients give exponent greater than $2$, and the coefficient
$\tfrac12L$ gives both a rational and a Liouville sum.

Fix $0<\theta<1/3$, and put
\begin{equation}\label{eq:sharpweights}
m_j=\lfloor\theta2^j\rfloor,\qquad
\lambda_j=N^{(-1)^jm_j},\qquad
e_j=2^j-(-1)^jm_j.
\end{equation}
The weights are positive rationals, with
\begin{equation}\label{eq:sharpheight}
\log H(\lambda_j)=m_jL=\theta L2^j+O(1).
\end{equation}
The associated sum has the exact digit expansion
\[
x=\sum_{j\ge0}N^{-e_j}.
\]
The exponents are strictly increasing: for odd $j$,
\[
e_{j+1}-e_j=2^j-m_j-m_{j+1}\ge(1-3\theta)2^j>0,
\]
and for even $j$ their difference is $2^j+m_j+m_{j+1}>0$.
The gaps tend to infinity, so the expansion is not eventually periodic
and $x$ is irrational.

Let $p_j/q_j$ be the truncation through index $j$, in lowest terms.
Its denominator is exactly $q_j=N^{e_j}$, because its numerator is
$1$ modulo $N$ and hence is coprime to every prime dividing $N$.
The increasing integer exponents give
\begin{equation}\label{eq:sharptail}
0<x-\frac{p_j}{q_j}
 \le\frac{N^{-e_{j+1}}}{1-N^{-1}}\le2N^{-e_{j+1}}.
\end{equation}
Along even indices,
\[
\frac{e_{j+1}}{e_j}\longrightarrow\frac{2(1+\theta)}{1-\theta}>2.
\]
For every $\nu<2(1+\theta)/(1-\theta)$, \eqref{eq:sharptail} is less than
$q_j^{-\nu}$ at all sufficiently large even indices. Consequently
\begin{equation}\label{eq:sharpmu}
\mu(x)\ge\frac{2(1+\theta)}{1-\theta}>2.
\end{equation}
Equivalently, for each $0<\kappa<L/3$ this gives
\[
\log H(\lambda_j)\sim \kappa2^j,\qquad
\mu(x)\ge\frac{2(L+\kappa)}{L-\kappa}.
\]
Given any $C_*>0$, choose $0<\kappa<\min(C_*,L/3)$.
Then the same weights satisfy $\log H(\lambda_j)\le C_*2^j$ for
all $j$, but their sum has exponent greater than $2$.

A rational sum is possible with positive integer weights. Set
\begin{equation}\label{eq:rationalsharp}
\lambda_0=N-1,\qquad
\lambda_j=N^{2^{j-1}}-1\quad(j\ge1).
\end{equation}
For $j\ge1$,
\[
\lambda_jN^{-2^j}=N^{-2^{j-1}}-N^{-2^j},
\]
so the partial sums telescope:
\[
\sum_{j=0}^{m}\lambda_jN^{-2^j}=1-N^{-2^m}.
\]
The infinite sum is therefore $1$, whereas
\begin{equation}\label{eq:halfheight}
\log H(\lambda_j)\sim\frac12L\,2^j.
\end{equation}

The same coefficient also allows a Liouville number. Keep $\lambda_0=N-1$, and replace the later weights by
\[
\widetilde\lambda_j=N^{2^{j-1}}-1+
\boldsymbol1_{\{j=r^2\text{ for some }r\ge1\}}\quad(j\ge1).
\]
These are again positive integers and satisfy \eqref{eq:halfheight}.
Their sum is
\[
\widetilde x=1+\xi_N,\qquad \xi_N=\sum_{r\ge1}N^{-2^{r^2}}.
\]
Its digit expansion is not eventually periodic, so $\xi_N$ is irrational.
The truncation through $r=R$ has reduced denominator
$q_R=N^{2^{R^2}}$, since its numerator is $1$ modulo $N$, and
\[
0<\xi_N-\sum_{r=1}^{R}N^{-2^{r^2}}
 \le2N^{-2^{(R+1)^2}}=2q_R^{-2^{2R+1}}.
\]
The same bound holds after adding $1$. Since $2^{2R+1}\to\infty$,
this proves $\mu(\widetilde x)=\infty$.

The height condition is not necessary for an individual series. For example, given an integer $b\ge2$,
the weights $\lambda_j=b^{-2^j}$ have
$\log H(\lambda_j)=(\log b)2^j$, but their sum equals
$\sum_{j\ge0}(Nb)^{-2^j}$. Theorem~\ref{thm:main}, applied with
base $Nb$ and constant weights $1$, gives exponent $2$ for that sum.


\section*{Use of AI tools}

The author devised the main argument. Codex (OpenAI, GPT-6 Astra) carried
out exact checks of the Hankel identities and searched the literature.
Claude Code (Anthropic, Claude Opus 5.5) reviewed the arguments and
references. The author takes responsibility for the proof and references.

\begin{thebibliography}{9}

\bibitem{AB}
M.~Amou and Y.~Bugeaud, \emph{Exponents of Diophantine approximation
and expansions in integer bases}, J. Lond. Math. Soc. (2)
\textbf{81} (2010), 297--316.

\bibitem{BHWY}
Y.~Bugeaud, G.-N.~Han, Z.-Y.~Wen, and J.-Y.~Yao,
\emph{Hankel determinants, Pad\'e approximations, and irrationality
exponents}, Int. Math. Res. Not. IMRN \textbf{2016} (2016), no.~5,
1467--1496, \url{https://doi.org/10.1093/imrn/rnv185}.

\bibitem{Du01}
D.~Duverney, \emph{Irrationality of fast converging series of rational
numbers}, J. Math. Sci. Univ. Tokyo \textbf{8} (2001), 275--316.

\bibitem{Du24}
D.~Duverney, \emph{Irrationality of certain fast converging series and
infinite products}, Hokkaido Math. J. \textbf{53} (2024), no.~2,
377--394, \url{https://doi.org/10.14492/hokmj/2022-679}.

\bibitem{DKS}
D.~Duverney, T.~Kurosawa, and I.~Shiokawa,
\emph{Irrationality exponents of certain fast converging series of
rational numbers}, Tsukuba J. Math. \textbf{44} (2020), no.~2,
235--250, \url{https://doi.org/10.21099/tkbjm/20204402235}.

\bibitem{Ni90}
K.~Nishioka, \emph{New approach in Mahler's method},
J. Reine Angew. Math. \textbf{407} (1990), 202--219.

\bibitem{Ra}
A.~Rajchert, \emph{On the irrationality exponents of Mahler numbers},
arXiv:2411.10733 (2024).

\bibitem{Sh}
J.~O.~Shallit, \emph{Simple continued fractions for some irrational
numbers}, J. Number Theory \textbf{11} (1979), no.~2, 209--217,
doi:10.1016/0022-314X(79)90040-4.

\bibitem{Su}
L.-H.~Su, \emph{Binary-valuation identities and lacunary continued fractions
in integer sequences}, Zenodo record 23676031 (2026),
\url{https://doi.org/10.5281/zenodo.23676031}.

\end{thebibliography}

\end{document}
